\section*{2 Preliminaries}
For integers $a,b$, we use $[a,b]$ to denote the set $\{a,\ldots,b\}$ and $[n]$ as a shorthand for $[1,n]$. Given a set $S\subseteq[n]$, when the ground set $[n]$ is clear from context, we denote $\rho(S)=|S|/n$, and its $\pm1$ variant as $\bias(S)=1-2\rho(S)$. For $z\in\mathbb F_2^n$, we similarly denote $\bias(z)$ as the bias of its characteristic set, i.e., $\mathbb E_{i\in[n]}[(-1)^{z_i}]$. We use boldface letters to denote random variables, except for $\mathbf1\in\mathbb R^n$, which we use for the all-ones vector. Also, when bounding running time, by writing $g(n)=\exp(f(n))$ we mean that $g(n)\leq2^{c\cdot f(n)}$ for some universal constant $c>0$.
\paragraph{Definition5 (discretizable distribution).} For $M\in\mathbb N$, We say that a distribution $\mathbf W$ is $M$-discretizable if it satisfies either of the following two equivalent properties.
\begin{enumerate}
\item For any $x$ in the support of $\mathbf W$, $\Pr_{\mathbf x\sim\mathbf W}[\mathbf x=x]=i/M$ for some $i\in\mathbb N$.
\item Let $U_M$ be the uniform distribution over $[M]$. Then, there is some function $f$ mapping $[M]$ to the support of $\mathbf W$ for which $f(i)$, where $i\sim U_M$, has the same distribution as a sample from $\mathbf W$.
\end{enumerate}
We say that $\mathbf W$ is computable in (deterministic or probabilistic) time $T$ if $f$ above is computable in time $T$. In particular, $\mathbf W$ is explicit if it is computable in deterministic time $\operatorname{poly}(M)$, and fully explicit if it is computable in deterministic time $\operatorname{poly}(\log M)$.
\paragraph{Definition6 (homogeneous distribution).} For $n,t\in\mathbb N$, we say that a distribution $\mathbf W$ over $[n]^t$ is homogeneous if its restriction to any coordinate is uniform over $[n]$, i.e., if for any $i\in[t]$ and $a\in[n]$ it holds that $\Pr_{\mathbf w\sim\mathbf W}[\mathbf w_i=a]=1/n$.
For any domain $\mathcal X$ and two distributions $\mathbf D,\mathbf D'$ over $\mathcal X$ we define the total variation distance of $\mathbf D$ and $\mathbf D'$ in terms of the optimal test distinguishing the distributions, i.e.,
\[d_{\mathrm{TV}}(\mathbf D,\mathbf D')\coloneqq\sup_{T:\mathcal X\to[0,1]}\{\mathbb E_{\mathbf x\sim\mathbf D}[T(\mathbf x)]-\mathbb E_{\mathbf x\sim\mathbf D'}[T(\mathbf x')]\}.\]
For a matrix $A\in\mathbb R^{n\times n}$, we denote by $\|A\|_{\mathrm{op}}$ its operator norm $\|A\|_{\mathrm{op}}=\max_{x\neq0}\|Ax\|_2/\|x\|_2$, which is also the maximum of $x^{\dagger}Ay$ over all norm-1 vectors $x,y\in\mathbb R^n$.
\paragraph{Error correcting codes.}
A binary error correcting code of message length $k$ and blocklength $n$ is a mapping $C:\mathbb F_2^k\to\mathbb F_2^n$, which we will often identify with its image $\operatorname{Im}(C)\subseteq\mathbb F_2^n$. The rate of $C$ is $k/n$, and its relative distance is $\Delta(C)=\frac1n\min_{z\neq z'}\Delta(z,z')$ for $z,z'\in C$, and $\Delta(z,z')=|\{i\in[n]:z_i\neq z'_i\}|$ being the Hamming distance. The Hamming ball of (relative) radius $\beta$ centered at $z$ is the set $B(z,\beta)=\{z'\in\mathbb F_2^n:\Delta(z,z')/n\leq\beta\}$.
We denote $\bias(C)$ as the maximal bias, in absolute value, of every nonzero $z\in C$. Thus, $\bias(C)\leq\varepsilon$ if the Hamming weight of any nonzero codeword is in $[(1-\varepsilon)/2,(1+\varepsilon)/2]$.
\paragraph{Definition7 (balanced codes).} A linear binary error correcting code is $\varepsilon$-balanced if $\bias(C)\leq\varepsilon$. In particular, an $\varepsilon$-balanced code $C$ has distance at least $(1-\varepsilon)/2$.
\paragraph{Unique and list decoding.}
We say that $C$ is (combinatorially) $(\beta,L)$ list decodable if for every $z\in\mathbb F_2^n$, $|C\cap B(z,\beta)|\leq L$. The Johnson bound tells us that any $\varepsilon$-balanced code is $(1/2-\sqrt\varepsilon,L)$ list decodable for $L=O(1/\varepsilon)$. The algorithmic list decoding problem aims at finding the list $L_{C,\beta}(z)\coloneqq C\cap B(z,\beta)$. When $\beta\leq(1-\varepsilon)/4$ and $C$ is $\varepsilon$-balanced, we know that $L_{C,\beta}$ always contains at most one codeword, which corresponds to the unique decoding problem.
\subsection*{2.1 Parity samplers and direct sum codes}
We will be interested in constructing a binary linear code $C\subseteq\mathbb F_2^{\bar n}$ with small bias by amplifying the (moderate) bias of some base code $C_0:\mathbb F_2^k\to\mathbb F_2^n$. One natural way to do so is by XORing $t$-tuples of $C_0$ according to some distribution $\mathbf W\sim[n]^t$.
\paragraph{Definition8 (direct sum codes).} For $t,n,\bar n\in\mathbb N$, let $\mathbf W\sim[n]^t$ be an $\bar n$-discretizable distribution equipped with a corresponding mapping function $f_{\mathbf W}:[\bar n]\to[n]^t$. For $z\in\mathbb F_2^n$, we let $\dsum_{\mathbf W}(z)\in\mathbb F_2^{\bar n}$ be such that
\[\dsum_{\mathbf W}(z)[\ell]=\sum_{i=1}^t z[f_{\mathbf W}(\ell)_i]\]
where the addition is taken over $\mathbb F_2$. Given a code $C_0\subseteq\mathbb F_2^n$, the direct sum lift of $C_0$ according to $\mathbf W$ is the code $\dsum_{\mathbf W}(C_0)=\{\dsum_{\mathbf W}(z):z\in C_0\}\subseteq\mathbb F_2^{\bar n}$.
\paragraph{Definition9 (parity sampler).} For $t,n\in\mathbb N$, and $0\leq\varepsilon<\varepsilon_0\leq1$, we say that $\mathbf W\sim[n]^t$ is an $(\varepsilon_0,\varepsilon)$ parity sampler if for every $z\in\mathbb F_2^n$ with $|\bias(z)|\leq\varepsilon_0$ it holds that $|\bias(\dsum_{\mathbf W}(z))|\leq\varepsilon$.
Clearly, if $C_0:\mathbb F_2^k\to\mathbb F_2^n$ is $\varepsilon_0$-balanced and $\mathbf W\sim[n]^t$ is an $\bar n$-discretizable $(\varepsilon_0,\varepsilon)$ parity sampler, the lifted code $C=\dsum_{\mathbf W}(C_0)$ is $\varepsilon$-balanced with rate $k/\bar n$.
\section*{3 Expanding 3-Uniform Hypergraphs}
Our construction uses a family of expanding $d$-regular3-uniform hypergraphs.
\paragraph{Definition10 ($d$-regular3-uniform hypergraph).} A3-uniform hypergraph consists of a set of vertices, $V$, and hyperedges $E\subseteq V^3$. The hypergraph $H=(V,E)$ is $d$-regular if, for each $v\in V$ and $j\in[3]$, the number of hyperedges $(w_1,w_2,w_3)\in E$ for which $v=w_j$ is $d$.
We will set $d$ carefully in our construction, but for now, it can be thought of as an arbitrary constant. For the remainder of this section, we will use ``hypergraph'' as shorthand for $d$-regular3-uniform hypergraph. There are various notions of expansion for hypergraphs and they are not all equivalent. In this work, we consider two such notions.
\paragraph{Definition11 ($\lambda$-mixing hypergraph).} A $d$-regular hypergraph $H=(V,E)$ on $n$ vertices is $\lambda$-mixing if, for any $S_1,S_2,S_3\subseteq V$,
\[\left|\frac{E(S_1,S_2,S_3)}d-\frac{|S_1|\cdot|S_2|\cdot|S_3|}{n^2}\right|\leq\lambda\sqrt{|S_1|\cdot|S_3|},\tag{2}\]
where $E(S_1,S_2,S_3)$ is the number of hyperedges $(w_1,w_2,w_3)\in E$ where $w_j\in S_j$ for $j=1,2,3$.
Note that the right-hand side of the above definition does not depend on the size of $S_2$.\footnote{Some works consider the stronger requirement of $\sqrt{|S_{\sigma(1)}|\cdot|S_{\sigma(2)}|}$ instead of $\sqrt{|S_1|\cdot|S_3|}$, for $S_{\sigma(1)}$ and $S_{\sigma(2)}$ being the two smallest sets.} That parallels the definition below, in which we use $\|y\|_{\infty}$ instead of $\|y\|_2$.
\paragraph{Definition12 ($\lambda$-spectral hypergraph).} A $d$-regular hypergraph $H=(V,E)$ on $n$ vertices is $\lambda$-spectral if, for any $x,y,z\in\mathbb R^n$,
\begin{multline*}
\left|\frac1d\sum_{(i,j,k)\in E}x_iy_jz_k-\frac1{n^2}\sum_{i\in V}x_i\sum_{j\in V}y_j\sum_{k\in V}z_k\right|\\
\leq\lambda\|x\|_2\|y\|_{\infty}\|z\|_2.\tag{3}
\end{multline*}
We will only care about cases where $y\in\{\pm1\}^n$.\footnote{In fact, for any fixed choice of $x$ and $z$, the left-hand side of Equation(3) is linear in $y$. Therefore, it is maximized for some $y\in\{\pm1\}^n$ and so in general it is sufficient to consider only such $y$.} In those cases, we have $\|y\|_{\infty}=1$ and $\|y\|_2=\sqrt n$. It is thus tempting to replace the $\|y\|_{\infty}$ in the right-hand side of Equation(3) with $\|y\|_2/\sqrt n$, as $\ell_2$ norms are often easier to work with than $\ell_{\infty}$ norms. Unfortunately, no ``good'' $\lambda$-spectral hypergraphs would exist with that modification: Whenever $n\gg d^2$, it is straightforward to bound $\lambda=\Omega(\sqrt n/d)$. For our definition, as we shall soon see, it is possible to achieve $\lambda\approx1/\sqrt d$.
Similar to standard graphs, spectral expansion implies mixing.
\paragraph{Proposition13 (spectral implies mixing).} For any $\lambda>0$, if $H$ is a $\lambda$-spectral hypergraph, it is also a $\lambda$-mixing hypergraph.
\begin{proof}
For any $S_1,S_2,S_3\subseteq V$, let $x_i=1[i\in S_1]$, $y_i=1[i\in S_2]$, and $z_i=1[i\in S_3]$. Then, $E(S_1,S_2,S_3)=\sum_{(i,j,k)\in E}x_iy_jz_k$. Equation(2) follows directly from Equation(3).
\end{proof}
By applying the converse to the expander mixing lemma for ordinary graphs[7], we can show that for symmetric hypergraphs, mixing implies spectral expansion with only a minor quantitative gap.
\paragraph{Definition14 (symmetric3-uniform hypergraph).} We say a3-uniform hypergraph $H=(V,E)$ is symmetric if, for any edge $e=(v_1,v_2,v_3)\in E$, the edge $(v_3,v_2,v_1)$ is also in $E$.
